\renewcommand{\thesection}{\arabic{section}}\renewcommand{\theHsection}{chapter.\arabic{section}}\setcounter{section}{3}\section{质点碰撞}
\noindent\begin{minipage}{\textwidth}\ProblemHeading{例题}{4.1（改写）}
一批速率均为 $V>0$ 的母粒子各自分裂成两个粒子。选定其中一种质量为 $m$ 的产物，它在母粒子的质心参考系中的速率恒为 $v_0>0$，飞出方向各向同性。求这一种产物在实验室参考系中的动能概率密度及动能范围。
\end{minipage}\par\vfill
\noindent\begin{minipage}{\textwidth}\ProblemHeading{例题}{4.2（改写）}
质量为 $M$ 的母粒子分裂成多于两个的粒子。在其质心参考系中，分裂放出的能量 $\varepsilon$ 全部成为各产物的动能。已知其中一个产物的质量为 $m_1$，求它的动能上限，并判断其余产物在什么情况下让它达到上限。
\end{minipage}\par\vfill
\newpage
\noindent\begin{minipage}{\textwidth}\ProblemHeading{习题}{4.3}
一个粒子以速率 $V>0$ 在实验室参考系中运动，并在无外力作用下分裂成质量分别为 $m_1$ 和 $m_2$ 的两个质点，分裂释放的能量为 $\varepsilon$。记两质点的速度方向与原运动方向所成的角分别为 $\theta_1$ 和 $\theta_2$。求 $\theta_1$ 和 $\theta_2$ 之间的关系。
\end{minipage}\par\vfill
\noindent\begin{minipage}{\textwidth}\ProblemHeading{习题}{4.4}
一个以速率 $V>0$ 在实验室参考系中运动的粒子分裂成两个质点。其中一种产物相对于质心的速率为 $v_0$，且飞出方向在质心参考系中各向同性。求该产物在实验室参考系中相对于原粒子运动方向的飞出角分布。
\end{minipage}\par\vfill
\newpage
\noindent\begin{minipage}{\textwidth}\ProblemHeading{习题}{4.5}
一个粒子以速率 $V$ 在实验室参考系中运动，随后分裂成两个质点。设两质点在质心参考系中的速率分别为 $v_{10}$ 和 $v_{20}$。求两质点在实验室参考系中飞出方向之间夹角的取值范围。
\end{minipage}\par\vfill
\noindent\begin{minipage}{\textwidth}\ProblemHeading{例题}{4.6（改写）}
质量为 $m_1$ 的质点以速率 $v$ 与初始静止、质量为 $m_2$ 的质点发生弹性碰撞。
\SubQ{1}{若两质点发生正碰，求它们碰后的速度。}
\SubQ{2}{在所有碰撞中，原来静止的质点至多能获得多少动能？}
\SubQ{3}{若两质点质量相等，且碰后速度均非零，证明两飞出方向垂直。}
\end{minipage}\par\vfill
\newpage
\noindent\begin{minipage}{\textwidth}\ProblemHeading{习题}{4.7}
质量为 $m_1$ 的质点以速率 $v$ 与一个原来静止、质量为 $m_2$ 的质点发生弹性碰撞。用实验室参考系中各质点碰后速度方向与入射方向的夹角，分别表示两个质点碰后的速率。
\end{minipage}\par\vfill
\noindent\begin{minipage}{\textwidth}\ProblemHeading{习题}{4.8}
求质点在半径为 $a$ 的刚性球上发生弹性散射的等效截面。质点与球的相互作用势能在 $r<a$ 时为 $U=\infty$，在 $r>a$ 时为 $U=0$，其中 $r$ 为质点到球心的距离。
\end{minipage}\par\vfill
\newpage
\noindent\begin{minipage}{\textwidth}\ProblemHeading{习题}{4.9}
质量为 $m_1$ 的质点以无穷远处的速率 $v_\infty$ 入射，在初始静止、质量为 $m_2$、半径为 $a$ 的刚性球上发生弹性散射。把等效截面表示为入射质点损失的动能 $\varepsilon$ 的函数。
\end{minipage}\par\vfill
\noindent\begin{minipage}{\textwidth}\ProblemHeading{习题}{4.10}
质点在势能 $U(r)\propto r^{-n}$（$n>0$）的有心力场中散射。在散射角固定时，等效截面怎样依赖于质点在无穷远处的速率 $v_\infty$？
\end{minipage}\par\vfill
\newpage
\noindent\begin{minipage}{\textwidth}\ProblemHeading{习题}{4.11}
质点从无穷远处以速率 $v_\infty$ 入射，在有心力场 $U(r)=-\alpha/r^2$（$\alpha>0$）中运动。求质点落向场中心的等效截面。
\end{minipage}\par\vfill
\noindent\begin{minipage}{\textwidth}\ProblemHeading{习题}{4.12}
质点从无穷远处以速率 $v_\infty$ 入射，在有心力场 $U(r)=-\alpha/r^n$（$n>2$，$\alpha>0$）中运动。求质点落向场中心的等效截面。
\end{minipage}\par\vfill
\newpage
\noindent\begin{minipage}{\textwidth}\ProblemHeading{习题}{4.13}
质量为 $m_1$ 的质点从无穷远处入射，和质量为 $m_2$、半径为 $R$ 的球体之间仅有牛顿引力。两者在无穷远处的相对速率为 $v_\infty$。求质点落到球体表面的等效截面。
\end{minipage}\par\vfill
\noindent\begin{minipage}{\textwidth}\ProblemHeading{习题}{4.14（校订）}
在给定的入射能量 $E>0$ 下，已知质点在有心力场中的微分散射截面 $\mathrm d\sigma/\mathrm d\chi$ 关于散射角 $\chi$ 的关系。设势能 $U(r)$ 随到中心的距离 $r$ 单调减小，并满足 $U(0)>E$、$U(\infty)=0$。求在该入射能量下可由散射实验探测的区域内的势能 $U(r)$。
\end{minipage}\par\vfill
\newpage
\noindent\begin{minipage}{\textwidth}\ProblemHeading{例题}{4.15（改写）}
质量为 $m$ 的质点在势能 $U(r)=\alpha/r$（$\alpha\ne0$）的固定有心力场中散射，无穷远处的入射速率为 $v_\infty$。
\SubQ{1}{求瞄准距离 $\rho$ 与散射角 $\chi$ 的关系。}
\SubQ{2}{求按出射方向的单位立体角计的微分散射截面，并判断其是否依赖于 $\alpha$ 的符号。}
\end{minipage}\par\vfill
\noindent\begin{minipage}{\textwidth}\ProblemHeading{例题}{4.16（改写）}
质量为 $m_1$ 的质点从无穷远处以速率 $v_\infty$ 入射，目标质点 $m_2$ 最初静止。两质点的相互作用势能为 $U(r)=\alpha/r$，其中 $r$ 为粒子间距且 $\alpha\ne0$。求下列情形在实验室参考系中相对于出射立体角的微分散射截面。
\SubQ{1}{当 $m_2\gg m_1$ 时，求入射质点的散射截面。}
\SubQ{2}{当 $m_1=m_2$ 且能区分两质点时，求入射质点的散射截面。}
\SubQ{3}{当两质点完全相同、实验只能记录出射粒子时，求按检测到的任一出射粒子计的截面。}
\end{minipage}\par\vfill
\newpage
\noindent\begin{minipage}{\textwidth}\ProblemHeading{例题}{4.17（改写）}
一个质量为 $m_1$ 的入射质点以远处速率 $v_\infty$ 飞向最初静止的质点 $m_2$。两质点的相互作用势能为 $U(r)=\alpha/r$（$\alpha\ne0$），其中 $r$ 为粒子间距。求在实验室参考系中以入射质点的动能损失 $\varepsilon$ 为变量的微分散射截面，以及 $\varepsilon$ 的取值范围。
\end{minipage}\par\vfill
\noindent\begin{minipage}{\textwidth}\ProblemHeading{习题}{4.18}
质量为 $m$ 的质点从无穷远处以速率 $v_\infty$ 入射，在固定有心力场 $U(r)=\alpha/r^2$（$\alpha>0$）中散射。求等效散射截面。
\end{minipage}\par\vfill
\newpage
\noindent\begin{minipage}{\textwidth}\ProblemHeading{习题}{4.19}
质量为 $m$ 的质点以无穷远处的速率 $v_\infty$ 入射，受到半径为 $a$、深度为 $U_0>0$ 的固定球形势阱散射。势能在 $r>a$ 时为 $U(r)=0$，在 $r<a$ 时为 $U(r)=-U_0$。求等效散射截面。
\end{minipage}\par\vfill
\noindent\begin{minipage}{\textwidth}\ProblemHeading{习题}{4.20（校订）}
质量为 $m_1$ 的质点以相对速率 $v_\infty$ 从无穷远处入射，目标质点 $m_2$ 最初静止，两者的相互作用势能为 $U(r)$，且 $U(\infty)=0$。设 $m$ 为折合质量，$\rho$ 为瞄准距离，$r_{\min}$ 为最近距离，质心参考系中的散射角为 $\chi=|\pi-2\varphi_0|$，其中
\[
\varphi_0=\int_{r_{\min}}^\infty
\frac{\rho\,\mathrm dr}{r^2\sqrt{1-\rho^2/r^2-2U(r)/(m v_\infty^2)}}.
\]
从这一积分表达式推导实验室参考系中的小偏转角公式
\[
\theta_1\simeq-\frac{2\rho}{m_1v_\infty^2}
\int_\rho^\infty\frac{U'(r)}{\sqrt{r^2-\rho^2}}\,\mathrm dr,
\]
其中 $\theta_1$ 取带符号的偏转角，朝远离散射中心的一侧偏转为正。
\end{minipage}\par\vfill
\newpage
\noindent\begin{minipage}{\textwidth}\ProblemHeading{习题}{4.21}
质量为 $m_1$ 的质点以速率 $v_\infty$ 入射，和初始静止的另一个质点发生小角度散射。已知两质点之间的势能为 $U(r)=\alpha/r^n$（$n>0$）。求实验室参考系中的等效散射截面。
\end{minipage}\par\vfill
